% ED26-BM-OHTR2012-1: actual published 2012 chapter1, pp.1--8.
\chapter{Thinking and Reasoning: A Source-Specific Field Comparison}
\label{sec:p11-thinking-and-reasoning-comparison}

Holyoak and Morrison's \emph{Thinking and Reasoning: A Reader's Guide}
organizes a cognitive research field around the transformation of mental
representations, overlapping reasoning tasks and different levels of
explanation. It begins with ordinary uses of thinking, proposes a provisional
scientific definition, sketches a historical development, and explains the
seven-part organization and several cross-cutting reading routes of the
2012 Oxford handbook. These are the guide's particular argument and choices,
not a universal partition of cognition or the organization of this reference
\parencite[pp.~1--8]{HolyoakMorrison2012ReadersGuide}.
\index{thinking!ordinary and scientific uses}\index{cognitive explanation!levels}
\index{reasoning!research traditions}

\section{Ordinary thinking does not have one fixed scientific meaning}
\label{sec:p11-ohtr1-ordinary}

The authors' opening invokes Descartes and the connection between thinking
and personal identity. They suggest that preserved intelligent action can
matter more to an ordinary assessment of mental competence than preserved
language alone. Their contrasts with sensory loss, language loss and serious
brain impairment motivate the inquiry; they are not clinical assessments,
legal capacity criteria or controlled comparisons establishing a hierarchy
of human worth \parencite[p.~1]{HolyoakMorrison2012ReadersGuide}.

The guide then makes the ambiguity concrete. Saying one thinks water is
needed for life expresses a belief; the same grammatical form can express a
less plausible belief about a political figure. Thinking of a solution
concerns constructing a way to attain a goal. An accusation that someone
failed to think before acting concerns foresight about possible consequences.
Asking what someone thinks of an option requests an evaluation; calling
genocide evil illustrates a moral judgment. Being lost in thought can instead
describe mental activity without a present practical problem.
The examples distinguish uses, not the truth or rationality of everything
that someone calls thinking \parencite[p.~1]{HolyoakMorrison2012ReadersGuide}.

\subsection{Kinds and parts are different questions}

The source's Figure~1.1, adapted from Rips and Conrad's 1989 study, summarizes
college students' judgments about mental-activity terms. One ordering asks
which activity is a kind of another; the other asks which is a part of
another. The reported relations can be re-expressed without treating the
figure as a map of brain mechanisms:

\begin{table}[htbp]
\centering\small
\begin{tabular}{p{.44\textwidth}p{.44\textwidth}}
Kinds relation & Parts relation\\ \hline
Planning is a kind of deciding. & Thinking is part of conceptualizing.\\
Deciding is a kind of reasoning. & Conceptualizing is part of remembering.\\
Reasoning is a kind of conceptualizing. & Remembering is part of reasoning.\\
Remembering is a kind of conceptualizing. & Remembering is part of reading.\\
Reading is a kind of conceptualizing. & Reasoning is part of deciding.\\
Conceptualizing is a kind of thinking. & Deciding is part of planning.
\end{tabular}
\caption{Independent verbal reconstruction of the two relation summaries in
the guide's inspected Figure~1.1. Corresponding rows across columns do not
assert paired measurements or causal relationships.}
\label{tab:p11-ohtr1-folk-relations}
\end{table}

This shows why a word can be broad in a category ordering while also naming
a component that people attribute to many activities. Planning can be viewed
as a kind of reasoning without implying that every reasoning episode is a
plan; treating thinking as part of remembering is not the assertion that
remembering is a subtype of thinking under exactly the same relation.
Conflating these questions makes apparent circularities look like evidence
against the subjects' meanings when two relations are actually being mixed.

\begin{proofboundary}{def:p11-ohtr1-folk-evidence}{A reported folk organization is not a cognitive ontology}
The guide identifies college students as the study population and judgments
of relations between terms as the task. It does not supply sample size,
recruitment, the full response instrument, analysis estimates or uncertainty
for the original study. The two summaries are therefore retained as the
guide's report, not independently quantified findings about all people,
languages or cultures. They also do not establish which operations cause
others or which representations a participant actually used.
\end{proofboundary}

The relevant reporting distinction is \cref{def:p19-empirical-question}:
a construct, response task and operational measure must be connected by
evidence, not identified because the same word occurs in all three.
\parencite[pp.~1--2, fig.~1.1]{HolyoakMorrison2012ReadersGuide}

\section{The provisional representational account and its boundaries}
\label{sec:p11-ohtr1-definition}

\begin{hbdefinition}{def:p11-ohtr1-thinking}{Systematic transformations of represented knowledge}
Under the editors' provisional account, thinking systematically transforms
internal descriptions carrying knowledge into descriptions of actual or
possible situations, often for a purpose. A mental representation is an
internal description usable in constructing other descriptions. The
transformations have constraints and can be characterized computationally;
thinking is not merely an arbitrary succession of associations.

Represented situations need not presently exist. They can concern an
imagined world, a future outcome or the thinker as an object of reflection.
Purpose is common but not mandatory: daydreaming and dreaming are not
excluded simply because a specified practical goal is absent.
\end{hbdefinition}

The definition is a proposal for organizing research, not a demonstration
that every process called thinking has one internal format or algorithm.
The authors expressly leave out several tempting requirements. Thinking
need not, by definition, require a human or a sentient being, although their
book chiefly examines biological human thinkers. They do not require all
thinking to be consciously accessible; some processes can be implicit.
They do not build rationality, optimality, desirability or intelligence into
the definition. A mistaken, biased or unwise transformation can still be
thinking. Disordered thought poses questions for the constraint-based account
rather than disappearing by a definition that counts only success
\parencite[pp.~1--2]{HolyoakMorrison2012ReadersGuide}.

The adjacent knowledge-representation treatment gives the needed explanatory
interface: \cref{def:p19-ckr-four-components} separates a representing
domain, a represented domain, their interpreting relations and processes
using them. An observer's ability to assign meaning to a mark does not show
that a modeled cognitive system has the same information or operation.
Similarly, a broad definition allowing nonhuman implementations does not
establish that any particular program thinks, is conscious or understands.

\begin{hbdefinition}{def:p11-ohtr1-tasks}{Inference, evaluation, choice and goal construction}
The guide distinguishes overlapping task perspectives:
\begin{itemize}
\item \emph{Reasoning} emphasizes drawing conclusions from initial
      information. Deductive validity, under the stated logic, guarantees
      truth of the conclusion from true premises by the argument form.
      Inductive support increases credibility without that guarantee.
\item \emph{Judgment} assesses an option's value or the probability of an
      outcome; \emph{decision making} selects among alternatives.
\item \emph{Problem solving} constructs a course of action to achieve a goal.
\end{itemize}
These perspectives are not disjoint boxes. A solution can require reasoning
about consequences, judging alternatives and deciding which action to take.
A logical inference task can itself be a problem with a goal.
\end{hbdefinition}

The source's note~2 says the deduction/induction distinction can blur in
psychological investigation. That does not make certainty and probabilistic
support equivalent: one experiment can mix interpretation, background
knowledge and several kinds of inference while formal standards remain
different. These guideposts orient the research question rather than dictate
which psychological process must occur in every task
\parencite[pp.~2--3, note~2 on p.~7]{HolyoakMorrison2012ReadersGuide}.

\section{The historical capsule: several traditions feed one field}
\label{sec:p11-ohtr1-history}

Holyoak and Morrison connect the philosophical study of thinking to modern
cognitive psychology, cognitive science and cognitive neuroscience. Their
capsule emphasizes selected developments, not an exhaustive priority history.
Hume and Kant frame the origins of causal knowledge; Hobbes's reckoning
metaphor motivates understanding reasoning through computation. The shared
word does not make seventeenth-century computation, a twentieth-century
computer simulation and a normative standard the same theory
\parencite[p.~3, note~3 on p.~7]{HolyoakMorrison2012ReadersGuide}.

In the early twentieth-century portion, the Gestalt tradition contributes
problem-solving and insight questions: Duncker and Wertheimer investigate
human problems, while K\"ohler examines great-ape problem solving. The guide
links Freud's legacy to unconscious thought, emotion and clinical inquiry;
Piaget and Vygotsky to development; Spearman to individual differences in
intelligence; and Luria to relationships between thinking and brain function.
These identify different phenomena and methods. Their inclusion is not
endorsement of every associated doctrine, an exclusive founder claim or
proof that a clinical theory has been experimentally established
\parencite[p.~3]{HolyoakMorrison2012ReadersGuide}.

Simon contributes bounded-resource alternatives to idealized economic
rationality: attention, memory and time constrain feasible operations.
Craik's 1943 model analogy supplies another bridge. Just as an engineer
can use a scale model of a bridge to explore stress, an internal model can
be run to anticipate a possible state. But a model's predictions depend on
what its construction preserves and what it leaves out; successful use of
one model is not proof that all thought consists of inspecting pictures.
The guide connects this tradition to Wason and Johnson-Laird's later
reasoning research \parencite[p.~3, note~4 on p.~7]{HolyoakMorrison2012ReadersGuide}.

Newell and Simon's problem-solving program makes the computational account
operational: initial conditions and a goal supply inputs, operations change
states, and a sequence of operations can reduce the difference between
current and desired conditions. Computer simulations and production-system
models thereby become ways to specify and test proposed processes. A
program's ability to reach the goal does not establish that a person follows
its particular search path; efficiency, errors and intermediate behavior
can distinguish candidate explanations
\parencite[pp.~3--4]{HolyoakMorrison2012ReadersGuide}.

The guide's 1970s account includes Rosch's work on categories and prototypes,
and Tversky and Kahneman's heuristics-and-biases research. Its whale/bear
comparison illustrates unusual versus typical mammals in a psychological
category. Atypicality must not be mistaken for uncertain biological
membership. Heuristics are described as often quick and useful yet capable
of nonnormative judgments in some circumstances. No universal error rate,
participant estimate or current consensus follows from this historical
overview. Relational reasoning, connectionist models, neuroscience and
neurogenetics are further research streams in the authors' early-new-
millennium field snapshot \parencite[p.~4]{HolyoakMorrison2012ReadersGuide}.

\section{Marr's cash register separates goal, procedure and realization}
\label{sec:p11-ohtr1-marr}

\begin{hbdefinition}{def:p11-ohtr1-levels}{Three different explanatory questions}
The \emph{computational level}, in the guide's presentation of Marr, asks
what a system is to accomplish and what input--output mapping serves that
task. It is not merely a count of processor operations.
The \emph{representation-and-algorithm level} asks how inputs and outputs
are encoded and by what procedure the transformation is carried out.
The \emph{implementation level} asks how that representation and procedure
are physically realized. Answers at one level can constrain another without
being substitutes for it.
\end{hbdefinition}

\begin{hbexample}{def:p11-ohtr1-cash-register}{A purchase total does not identify its code or hardware}
The source's cash register is intended to calculate the amount owed for a
purchase under a specified pricing arrangement. Adding item prices explains
the desired mapping; invariance to the order of presenting items is one
relevant constraint. It is not, alone, a characterization of addition:
multiplication is order-independent too. The purchase-total specification
and its other assumptions establish which mapping is intended.

That task can be represented in decimal or binary and performed by different
addition procedures. It can be realized electronically, mechanically or
through a clerk's mental abacus. Correct totals alone do not identify the
code, the sequence of operations or the physical realization.
\end{hbexample}

Marr's warning in this account is against focusing prematurely on the concrete
mechanism while leaving the task unexplained. It is not a claim that neural
or behavioral evidence is dispensable. The guide's note~5 records Marr's
criticism that Newell and Simon gave too little attention to the computational
level in this sense. Thus two computational approaches can disagree about
the appropriate level of analysis rather than differ merely in software
implementation \parencite[p.~4, note~5 on p.~7]{HolyoakMorrison2012ReadersGuide}.

\section{The actual seven-part 2012 organization and current routes}
\label{sec:p11-ohtr1-organization}

The published guide organizes its handbook into seven parts. The chapter
numbers below belong to that Oxford volume; they are not part or chapter
numbers of \emph{The Methods of Reason}.

\paragraph{I. General approaches, chapters 2--7.}
The guide begins with normative logic, probability and rational choice,
then Bayesian inference, knowledge representation and computational models,
followed by cognitive neuroscience and neurogenetics. These place standards,
representations, processes and realizations alongside one another, not in a
single ascending list of methods that replace their predecessors. Current
interfaces include \cref{sec:p11-rationality-intro},
\cref{sec:p15-bayesian-updating-prediction},
\cref{sec:p19-cognitive-knowledge-representation} and
\cref{sec:p19-evidence-methods}; each answers a different question.

\paragraph{II. Inductive, deductive and abductive reasoning, chapters 8--15.}
The source includes dual-process accounts, mental models, similarity,
concepts/categories, causal learning, analogy, explanation and Bayesian
approaches to informal argument. Their placement shows that a logical
standard, a representational hypothesis and an explanation of an argument's
persuasiveness are not identical. The formal abduction route
\cref{sec:p16-finite-abduction} asks which hypotheses explain a specified
observation under declared constraints; structured argumentation
\cref{sec:p18-structured-argumentation} asks about arguments, challenges
and defeat. Neither title by itself supplies every cognitive account in
the source's part.

\paragraph{III. Judgment and decision making, chapters 16--20.}
Decision making, heuristics/biases, behavioral game theory and neuroeconomics,
moral judgment, and motivation/emotion bring evaluation and choice into the
field. Use \cref{sec:p17-finite-expected-utility} to evaluate a supplied
lottery preference model and \cref{sec:p17-finite-strategic-games} for
strategic deviations under a supplied game. Use an empirical evidence record
to ask what people did. An observed choice does not justify a normative
utility scale, and an equilibrium calculation is not evidence that people
play it.

\paragraph{IV. Problem solving, intelligence and creative thinking, chapters 21--25.}
The guide relates general problem solving to rationality/intelligence,
creativity, insight and the extreme individual differences described as
genius. The tasks include generating an idea, finding a solution and evaluating
reasoning, which need not share a single performance measure. The distinction
between capacity and observed performance in
\cref{def:p11-ri-capacity-performance}, and between construct and task
in \cref{def:p19-empirical-question}, prevents an intelligence score or
one successful solution from becoming a complete assessment of rationality.

\paragraph{V. Ontogeny, phylogeny, language and culture, chapters 26--29.}
Development, comparative cognition, language/thought and socially distributed
thinking challenge restriction to young adult English speakers. They call
for population and task comparisons rather than assuming that an unchanged
label denotes an invariant construct across ages, species or languages.
\Cref{def:p19-population-and-estimand,sec:p19-comparability-and-replication}
supplies the current measurement and transfer boundary; it does not settle
the substantive hypotheses about those populations.

\paragraph{VI. Modes of thinking, chapters 30--34.}
Mathematical cognition, visuospatial thinking, gesture, aging and thought
disorder distinguish variations in representation and process. The
visualization comparison in \cref{sec:p19-rationality-13-3} separates
pictorial detail from spatial organization and task performance from proof.
The guide's clinical and developmental topic descriptions are not diagnostic
criteria, treatment recommendations or evidence of one universal trajectory.

\paragraph{VII. Thinking in practice, chapters 35--40.}
Science, legal reasoning, medicine, business, music and learning to think
show why an application needs its own conditions. For example,
\cref{sec:p11-rationality-11-3} distinguishes legal interpretation and
authority from formal validity. Empirical transfer or training claims need
\cref{def:p19-evidence-record}, not a presumption that practice in one
domain improves all reasoning. Identifying these source topics is a map of
questions, not a claim that every application is reconstructed in this entry.
\parencite[pp.~4--6]{HolyoakMorrison2012ReadersGuide}

\subsection{The guide's reading groups cross the part boundaries}

The authors give five example chapter groupings rather than one mandatory
order. Their actual selections are:
\begin{table}[htbp]
\centering\small
\begin{tabular}{p{.29\textwidth}p{.60\textwidth}}
Reading purpose & Oxford chapter numbers\\ \hline
Introduction &1, 2, 3, 4, 8, 9, 10, 11, 12, 13, 14, 15, 16, 17, 21, 22, 23\\
Development &1, 4, 10, 11, 13, 14, 26, 27, 28, 30, 31, 32, 33, 40\\
Modeling &1, 3, 4, 5, 6, 9, 10, 11, 12, 13, 15, 18, 40\\
Applied thought &1, 35, 36, 37, 38, 39, 40\\
Differences in thought &1, 19, 20, 23, 24, 25, 26, 27, 28, 29, 32, 33, 34
\end{tabular}
\caption{The five suggested chapter routes in the actual 2012 guide,
re-expressed by purpose. These are source-volume reading selections, not
prerequisites or a prescribed curriculum for this reference.}
\label{tab:p11-ohtr1-reading-routes}
\end{table}

Modeling spans representation, Bayesian inference, causal/relational
reasoning, behavioral games and learning. Development spans language,
spatial/mathematical cognition, gesture, aging and transfer. Applied thought
groups practice domains, while differences in thought groups motivational,
creative, developmental, cultural and clinical questions. This is evidence
of the editors' cross-cutting organization, not evidence that taking one
course improves all those capacities
\parencite[pp.~6--7]{HolyoakMorrison2012ReadersGuide}.

\section{Comparison with the rationality framework and task lookup}
\label{sec:p11-ohtr1-comparison}

The Oxford guide's primary question is what thinking comprises and how its
research traditions can be organized. The Knauff--Spohn framework in
\cref{sec:p11-rationality-intro} instead foregrounds what is assessed as
rational, which standard applies and how normative and descriptive inquiries
connect. They overlap on representations, capacities and Marr's levels but
do not have identical organizing axes. The Oxford definition explicitly
allows irrational thinking; it therefore cannot itself supply a standard
of rational success. The 2021 process/output distinction in
\cref{sec:p11-ri-process-output} deliberately permits normative goal
specification, while its theoretical/practical and individual/social
distinctions ask further questions not answered by Oxford's seven parts.

\begin{hbmethod}{def:p11-ohtr1-lookup}{Locate the question before transferring a field label}
If the question concerns the meanings of mental-activity words, distinguish
kinds from parts and treat the reported judgments as evidence about those
meanings, not neural architecture. If it concerns a cognitive explanation,
separate task/mapping, representation/procedure and realization. If it
concerns a justified conclusion or choice, state the appropriate formal or
normative standard rather than inherit it from the word thinking. If it
concerns actual performance or transfer, specify task, population, measure,
comparison and uncertainty before drawing the empirical conclusion.
\end{hbmethod}

A successful use of the guide distinguishes a cash register's intended
total from the algorithm implementing it, or separates a folk category
judgment from an observed inference process. A failed use assumes that
because an activity counts as thinking it is rational, that a simulation
reaching the right answer uniquely explains human thought, or that a
source-volume reading route proves either cognitive improvement or the
completeness of this reference's coverage.
